
% Compteur de séance partagé entre les appels successifs de ce fichier (ce
% fichier est \input plusieurs fois depuis fiche-evaluation.tex : chaque appel
% doit reprendre la numérotation des séances là où le précédent s'est arrêté,
% au lieu de recommencer à S1 à chaque fois).
\makeatletter
\@ifundefined{c@seanceoffset}{%
  \newcounter{seanceoffset}
  \setcounter{seanceoffset}{0}
}{}
\makeatother
\providecommand{\seancecol}[1]{S\number\numexpr\value{seanceoffset}+#1\relax}

\textbf{Nom :} \makebox[5cm]{\hrulefill} \hfill \textbf{Prénom :} \makebox[5cm]{\hrulefill}
\\
\\
Légende : \textbf{A} = Acquis, \textbf{ECA} = En cours d'acquisition, \textbf{NT} = Non
travaillé.

\begin{center}
  \footnotesize
  \renewcommand{\arraystretch}{1.25}
  \begin{longtable}{|p{3.2cm}|p{5.3cm}|c|c|c|c|c|c|c|c|}

    \caption{Fiche de suivi et d'évaluation détaillée --- Plongeur Niveau 1 (PE20)}\\
    \hline
    \rowcolor{ffessmblue}
    \textcolor{white}{\textbf{Compétence}} & \textcolor{white}{\textbf{Sous-compétence}} &
    \textcolor{white}{\textbf{\seancecol{1}}} & \textcolor{white}{\textbf{\seancecol{2}}} &
    \textcolor{white}{\textbf{\seancecol{3}}} & \textcolor{white}{\textbf{\seancecol{4}}} &
    \textcolor{white}{\textbf{\seancecol{5}}} & \textcolor{white}{\textbf{\seancecol{6}}} &
    \textcolor{white}{\textbf{\seancecol{7}}} & \textcolor{white}{\textbf{\seancecol{8}}} \\
    \hline
    % Table tenant toujours sur une seule page : un unique en-tête suffit,
    % pas besoin de \endfirsthead distinct ni de pied de page "suite ...".
    \endhead

    \hline
    \endlastfoot

    % --- S'équiper et se déséquiper ---
    \multirow{3}{=}{S'équiper et se déséquiper} & Gréage et dégréage & & & & & & & & \\ \cline{2-10}
    & Capelage et décapelage & & & & & & & & \\ \cline{2-10}
    & Choix de son matériel personnel & & & & & & & & \\
    \hline

    % --- Se mettre à l'eau et sortir de l'eau ---
    \multirow{4}{=}{Se mettre à l'eau et sortir de l'eau} & Saut droit & & & & & & & & \\ \cline{2-10}
    & Bascule arrière & & & & & & & & \\ \cline{2-10}
    & Départ plage & & & & & & & & \\ \cline{2-10}
    & Sortir de l'eau & & & & & & & & \\
    \hline

    % --- S'immerger ---
    \multirow{2}{=}{Évoluer dans l'eau --- S'immerger} & Canard & & & & & & & & \\ \cline{2-10}
    & Phoque & & & & & & & & \\
    \hline

    % --- Se propulser ---
    \multirow{5}{=}{Évoluer dans l'eau --- Se propulser} & Palmage ventral en surface & & & & & & & & \\ \cline{2-10}
    & Palmage dorsal & & & & & & & & \\ \cline{2-10}
    & Palmage de sustentation & & & & & & & & \\ \cline{2-10}
    & Palmage en immersion & & & & & & & & \\ \cline{2-10}
    & Nage en capelé & & & & & & & & \\
    \hline

    % --- Se ventiler ---
    \multirow{4}{=}{Évoluer dans l'eau --- Se ventiler} & Ventilation en immersion & & & & & & & & \\ \cline{2-10}
    & Ventilation sur tuba et vidage du tuba & & & & & & & & \\ \cline{2-10}
    & Vidage du masque & & & & & & & & \\ \cline{2-10}
    & Lâcher et reprise d'embout & & & & & & & & \\
    \hline

    % --- S'équilibrer ---
    \multirow{2}{=}{Évoluer dans l'eau --- S'équilibrer} & Gestion du gilet de stabilisation & & & & & & & & \\ \cline{2-10}
    & Poumon ballast & & & & & & & & \\
    \hline

    % --- Respecter le milieu et l'environnement ---
    Respecter le milieu et l'environnement & Aisance aquatique & & & & & & & & \\
    \hline

    % --- Communiquer ---
    Communiquer & Exécution des signes conventionnels & & & & & & & & \\
    \hline

    % --- Retourner en surface ---
    \multirow{5}{=}{Retourner en surface} & Maîtrise de la vitesse de remontée & & & & & & & & \\ \cline{2-10}
    & Tenue d'un palier & & & & & & & & \\ \cline{2-10}
    & Tour d'horizon & & & & & & & & \\ \cline{2-10}
    & Gonflage du gilet en surface & & & & & & & & \\ \cline{2-10}
    & Remontée en expiration contrôlée & & & & & & & & \\
    \hline

    % --- Évoluer en sécurité ---
    \multirow{2}{=}{Évoluer en sécurité} & Application des procédures mises en œuvre par le GP & & & & & & & & \\ \cline{2-10}
    & Intervention en relais & & & & & & & & \\
    \hline

  \end{longtable}
\end{center}

L'évaluation des compétences se fait en \textbf{contrôle continu}. Le niveau est validé
lorsque toutes les compétences sont acquises. Les connaissances théoriques sont évaluées à
l'oral, lors des mises en situations pratiques. Il n'y a pas d'examen écrit. L'accent est mis
sur la prévention.

\vspace{0.5cm}
\textbf{Validé le :} \makebox[6cm]{\hrulefill} \hfill \\
\\
\textbf{Nom et signature du moniteur :} \makebox[6cm]{\hrulefill}
\vspace{0.5cm}

{\small\itshape Cette fiche n'a aucun caractère obligatoire, elle est modulable et il
appartient au moniteur de l'adapter pour sa propre utilisation.}

% Prépare la numérotation pour le prochain 
% Compteur de séance partagé entre les appels successifs de ce fichier (ce
% fichier est \input plusieurs fois depuis fiche-evaluation.tex : chaque appel
% doit reprendre la numérotation des séances là où le précédent s'est arrêté,
% au lieu de recommencer à S1 à chaque fois).
\makeatletter
\@ifundefined{c@seanceoffset}{%
  \newcounter{seanceoffset}
  \setcounter{seanceoffset}{0}
}{}
\makeatother
\providecommand{\seancecol}[1]{S\number\numexpr\value{seanceoffset}+#1\relax}

\textbf{Nom :} \makebox[5cm]{\hrulefill} \hfill \textbf{Prénom :} \makebox[5cm]{\hrulefill}
\\
\\
Légende : \textbf{A} = Acquis, \textbf{ECA} = En cours d'acquisition, \textbf{NT} = Non
travaillé.

\begin{center}
  \footnotesize
  \renewcommand{\arraystretch}{1.25}
  \begin{longtable}{|p{3.2cm}|p{5.3cm}|c|c|c|c|c|c|c|c|}

    \caption{Fiche de suivi et d'évaluation détaillée --- Plongeur Niveau 1 (PE20)}\\
    \hline
    \rowcolor{ffessmblue}
    \textcolor{white}{\textbf{Compétence}} & \textcolor{white}{\textbf{Sous-compétence}} &
    \textcolor{white}{\textbf{\seancecol{1}}} & \textcolor{white}{\textbf{\seancecol{2}}} &
    \textcolor{white}{\textbf{\seancecol{3}}} & \textcolor{white}{\textbf{\seancecol{4}}} &
    \textcolor{white}{\textbf{\seancecol{5}}} & \textcolor{white}{\textbf{\seancecol{6}}} &
    \textcolor{white}{\textbf{\seancecol{7}}} & \textcolor{white}{\textbf{\seancecol{8}}} \\
    \hline
    % Table tenant toujours sur une seule page : un unique en-tête suffit,
    % pas besoin de \endfirsthead distinct ni de pied de page "suite ...".
    \endhead

    \hline
    \endlastfoot

    % --- S'équiper et se déséquiper ---
    \multirow{3}{=}{S'équiper et se déséquiper} & Gréage et dégréage & & & & & & & & \\ \cline{2-10}
    & Capelage et décapelage & & & & & & & & \\ \cline{2-10}
    & Choix de son matériel personnel & & & & & & & & \\
    \hline

    % --- Se mettre à l'eau et sortir de l'eau ---
    \multirow{4}{=}{Se mettre à l'eau et sortir de l'eau} & Saut droit & & & & & & & & \\ \cline{2-10}
    & Bascule arrière & & & & & & & & \\ \cline{2-10}
    & Départ plage & & & & & & & & \\ \cline{2-10}
    & Sortir de l'eau & & & & & & & & \\
    \hline

    % --- S'immerger ---
    \multirow{2}{=}{Évoluer dans l'eau --- S'immerger} & Canard & & & & & & & & \\ \cline{2-10}
    & Phoque & & & & & & & & \\
    \hline

    % --- Se propulser ---
    \multirow{5}{=}{Évoluer dans l'eau --- Se propulser} & Palmage ventral en surface & & & & & & & & \\ \cline{2-10}
    & Palmage dorsal & & & & & & & & \\ \cline{2-10}
    & Palmage de sustentation & & & & & & & & \\ \cline{2-10}
    & Palmage en immersion & & & & & & & & \\ \cline{2-10}
    & Nage en capelé & & & & & & & & \\
    \hline

    % --- Se ventiler ---
    \multirow{4}{=}{Évoluer dans l'eau --- Se ventiler} & Ventilation en immersion & & & & & & & & \\ \cline{2-10}
    & Ventilation sur tuba et vidage du tuba & & & & & & & & \\ \cline{2-10}
    & Vidage du masque & & & & & & & & \\ \cline{2-10}
    & Lâcher et reprise d'embout & & & & & & & & \\
    \hline

    % --- S'équilibrer ---
    \multirow{2}{=}{Évoluer dans l'eau --- S'équilibrer} & Gestion du gilet de stabilisation & & & & & & & & \\ \cline{2-10}
    & Poumon ballast & & & & & & & & \\
    \hline

    % --- Respecter le milieu et l'environnement ---
    Respecter le milieu et l'environnement & Aisance aquatique & & & & & & & & \\
    \hline

    % --- Communiquer ---
    Communiquer & Exécution des signes conventionnels & & & & & & & & \\
    \hline

    % --- Retourner en surface ---
    \multirow{5}{=}{Retourner en surface} & Maîtrise de la vitesse de remontée & & & & & & & & \\ \cline{2-10}
    & Tenue d'un palier & & & & & & & & \\ \cline{2-10}
    & Tour d'horizon & & & & & & & & \\ \cline{2-10}
    & Gonflage du gilet en surface & & & & & & & & \\ \cline{2-10}
    & Remontée en expiration contrôlée & & & & & & & & \\
    \hline

    % --- Évoluer en sécurité ---
    \multirow{2}{=}{Évoluer en sécurité} & Application des procédures mises en œuvre par le GP & & & & & & & & \\ \cline{2-10}
    & Intervention en relais & & & & & & & & \\
    \hline

  \end{longtable}
\end{center}

L'évaluation des compétences se fait en \textbf{contrôle continu}. Le niveau est validé
lorsque toutes les compétences sont acquises. Les connaissances théoriques sont évaluées à
l'oral, lors des mises en situations pratiques. Il n'y a pas d'examen écrit. L'accent est mis
sur la prévention.

\vspace{0.5cm}
\textbf{Validé le :} \makebox[6cm]{\hrulefill} \hfill \\
\\
\textbf{Nom et signature du moniteur :} \makebox[6cm]{\hrulefill}
\vspace{0.5cm}

{\small\itshape Cette fiche n'a aucun caractère obligatoire, elle est modulable et il
appartient au moniteur de l'adapter pour sa propre utilisation.}

% Prépare la numérotation pour le prochain 
% Compteur de séance partagé entre les appels successifs de ce fichier (ce
% fichier est \input plusieurs fois depuis fiche-evaluation.tex : chaque appel
% doit reprendre la numérotation des séances là où le précédent s'est arrêté,
% au lieu de recommencer à S1 à chaque fois).
\makeatletter
\@ifundefined{c@seanceoffset}{%
  \newcounter{seanceoffset}
  \setcounter{seanceoffset}{0}
}{}
\makeatother
\providecommand{\seancecol}[1]{S\number\numexpr\value{seanceoffset}+#1\relax}

\textbf{Nom :} \makebox[5cm]{\hrulefill} \hfill \textbf{Prénom :} \makebox[5cm]{\hrulefill}
\\
\\
Légende : \textbf{A} = Acquis, \textbf{ECA} = En cours d'acquisition, \textbf{NT} = Non
travaillé.

\begin{center}
  \footnotesize
  \renewcommand{\arraystretch}{1.25}
  \begin{longtable}{|p{3.2cm}|p{5.3cm}|c|c|c|c|c|c|c|c|}

    \caption{Fiche de suivi et d'évaluation détaillée --- Plongeur Niveau 1 (PE20)}\\
    \hline
    \rowcolor{ffessmblue}
    \textcolor{white}{\textbf{Compétence}} & \textcolor{white}{\textbf{Sous-compétence}} &
    \textcolor{white}{\textbf{\seancecol{1}}} & \textcolor{white}{\textbf{\seancecol{2}}} &
    \textcolor{white}{\textbf{\seancecol{3}}} & \textcolor{white}{\textbf{\seancecol{4}}} &
    \textcolor{white}{\textbf{\seancecol{5}}} & \textcolor{white}{\textbf{\seancecol{6}}} &
    \textcolor{white}{\textbf{\seancecol{7}}} & \textcolor{white}{\textbf{\seancecol{8}}} \\
    \hline
    % Table tenant toujours sur une seule page : un unique en-tête suffit,
    % pas besoin de \endfirsthead distinct ni de pied de page "suite ...".
    \endhead

    \hline
    \endlastfoot

    % --- S'équiper et se déséquiper ---
    \multirow{3}{=}{S'équiper et se déséquiper} & Gréage et dégréage & & & & & & & & \\ \cline{2-10}
    & Capelage et décapelage & & & & & & & & \\ \cline{2-10}
    & Choix de son matériel personnel & & & & & & & & \\
    \hline

    % --- Se mettre à l'eau et sortir de l'eau ---
    \multirow{4}{=}{Se mettre à l'eau et sortir de l'eau} & Saut droit & & & & & & & & \\ \cline{2-10}
    & Bascule arrière & & & & & & & & \\ \cline{2-10}
    & Départ plage & & & & & & & & \\ \cline{2-10}
    & Sortir de l'eau & & & & & & & & \\
    \hline

    % --- S'immerger ---
    \multirow{2}{=}{Évoluer dans l'eau --- S'immerger} & Canard & & & & & & & & \\ \cline{2-10}
    & Phoque & & & & & & & & \\
    \hline

    % --- Se propulser ---
    \multirow{5}{=}{Évoluer dans l'eau --- Se propulser} & Palmage ventral en surface & & & & & & & & \\ \cline{2-10}
    & Palmage dorsal & & & & & & & & \\ \cline{2-10}
    & Palmage de sustentation & & & & & & & & \\ \cline{2-10}
    & Palmage en immersion & & & & & & & & \\ \cline{2-10}
    & Nage en capelé & & & & & & & & \\
    \hline

    % --- Se ventiler ---
    \multirow{4}{=}{Évoluer dans l'eau --- Se ventiler} & Ventilation en immersion & & & & & & & & \\ \cline{2-10}
    & Ventilation sur tuba et vidage du tuba & & & & & & & & \\ \cline{2-10}
    & Vidage du masque & & & & & & & & \\ \cline{2-10}
    & Lâcher et reprise d'embout & & & & & & & & \\
    \hline

    % --- S'équilibrer ---
    \multirow{2}{=}{Évoluer dans l'eau --- S'équilibrer} & Gestion du gilet de stabilisation & & & & & & & & \\ \cline{2-10}
    & Poumon ballast & & & & & & & & \\
    \hline

    % --- Respecter le milieu et l'environnement ---
    Respecter le milieu et l'environnement & Aisance aquatique & & & & & & & & \\
    \hline

    % --- Communiquer ---
    Communiquer & Exécution des signes conventionnels & & & & & & & & \\
    \hline

    % --- Retourner en surface ---
    \multirow{5}{=}{Retourner en surface} & Maîtrise de la vitesse de remontée & & & & & & & & \\ \cline{2-10}
    & Tenue d'un palier & & & & & & & & \\ \cline{2-10}
    & Tour d'horizon & & & & & & & & \\ \cline{2-10}
    & Gonflage du gilet en surface & & & & & & & & \\ \cline{2-10}
    & Remontée en expiration contrôlée & & & & & & & & \\
    \hline

    % --- Évoluer en sécurité ---
    \multirow{2}{=}{Évoluer en sécurité} & Application des procédures mises en œuvre par le GP & & & & & & & & \\ \cline{2-10}
    & Intervention en relais & & & & & & & & \\
    \hline

  \end{longtable}
\end{center}

L'évaluation des compétences se fait en \textbf{contrôle continu}. Le niveau est validé
lorsque toutes les compétences sont acquises. Les connaissances théoriques sont évaluées à
l'oral, lors des mises en situations pratiques. Il n'y a pas d'examen écrit. L'accent est mis
sur la prévention.

\vspace{0.5cm}
\textbf{Validé le :} \makebox[6cm]{\hrulefill} \hfill \\
\\
\textbf{Nom et signature du moniteur :} \makebox[6cm]{\hrulefill}
\vspace{0.5cm}

{\small\itshape Cette fiche n'a aucun caractère obligatoire, elle est modulable et il
appartient au moniteur de l'adapter pour sa propre utilisation.}

% Prépare la numérotation pour le prochain 
% Compteur de séance partagé entre les appels successifs de ce fichier (ce
% fichier est \input plusieurs fois depuis fiche-evaluation.tex : chaque appel
% doit reprendre la numérotation des séances là où le précédent s'est arrêté,
% au lieu de recommencer à S1 à chaque fois).
\makeatletter
\@ifundefined{c@seanceoffset}{%
  \newcounter{seanceoffset}
  \setcounter{seanceoffset}{0}
}{}
\makeatother
\providecommand{\seancecol}[1]{S\number\numexpr\value{seanceoffset}+#1\relax}

\textbf{Nom :} \makebox[5cm]{\hrulefill} \hfill \textbf{Prénom :} \makebox[5cm]{\hrulefill}
\\
\\
Légende : \textbf{A} = Acquis, \textbf{ECA} = En cours d'acquisition, \textbf{NT} = Non
travaillé.

\begin{center}
  \footnotesize
  \renewcommand{\arraystretch}{1.25}
  \begin{longtable}{|p{3.2cm}|p{5.3cm}|c|c|c|c|c|c|c|c|}

    \caption{Fiche de suivi et d'évaluation détaillée --- Plongeur Niveau 1 (PE20)}\\
    \hline
    \rowcolor{ffessmblue}
    \textcolor{white}{\textbf{Compétence}} & \textcolor{white}{\textbf{Sous-compétence}} &
    \textcolor{white}{\textbf{\seancecol{1}}} & \textcolor{white}{\textbf{\seancecol{2}}} &
    \textcolor{white}{\textbf{\seancecol{3}}} & \textcolor{white}{\textbf{\seancecol{4}}} &
    \textcolor{white}{\textbf{\seancecol{5}}} & \textcolor{white}{\textbf{\seancecol{6}}} &
    \textcolor{white}{\textbf{\seancecol{7}}} & \textcolor{white}{\textbf{\seancecol{8}}} \\
    \hline
    % Table tenant toujours sur une seule page : un unique en-tête suffit,
    % pas besoin de \endfirsthead distinct ni de pied de page "suite ...".
    \endhead

    \hline
    \endlastfoot

    % --- S'équiper et se déséquiper ---
    \multirow{3}{=}{S'équiper et se déséquiper} & Gréage et dégréage & & & & & & & & \\ \cline{2-10}
    & Capelage et décapelage & & & & & & & & \\ \cline{2-10}
    & Choix de son matériel personnel & & & & & & & & \\
    \hline

    % --- Se mettre à l'eau et sortir de l'eau ---
    \multirow{4}{=}{Se mettre à l'eau et sortir de l'eau} & Saut droit & & & & & & & & \\ \cline{2-10}
    & Bascule arrière & & & & & & & & \\ \cline{2-10}
    & Départ plage & & & & & & & & \\ \cline{2-10}
    & Sortir de l'eau & & & & & & & & \\
    \hline

    % --- S'immerger ---
    \multirow{2}{=}{Évoluer dans l'eau --- S'immerger} & Canard & & & & & & & & \\ \cline{2-10}
    & Phoque & & & & & & & & \\
    \hline

    % --- Se propulser ---
    \multirow{5}{=}{Évoluer dans l'eau --- Se propulser} & Palmage ventral en surface & & & & & & & & \\ \cline{2-10}
    & Palmage dorsal & & & & & & & & \\ \cline{2-10}
    & Palmage de sustentation & & & & & & & & \\ \cline{2-10}
    & Palmage en immersion & & & & & & & & \\ \cline{2-10}
    & Nage en capelé & & & & & & & & \\
    \hline

    % --- Se ventiler ---
    \multirow{4}{=}{Évoluer dans l'eau --- Se ventiler} & Ventilation en immersion & & & & & & & & \\ \cline{2-10}
    & Ventilation sur tuba et vidage du tuba & & & & & & & & \\ \cline{2-10}
    & Vidage du masque & & & & & & & & \\ \cline{2-10}
    & Lâcher et reprise d'embout & & & & & & & & \\
    \hline

    % --- S'équilibrer ---
    \multirow{2}{=}{Évoluer dans l'eau --- S'équilibrer} & Gestion du gilet de stabilisation & & & & & & & & \\ \cline{2-10}
    & Poumon ballast & & & & & & & & \\
    \hline

    % --- Respecter le milieu et l'environnement ---
    Respecter le milieu et l'environnement & Aisance aquatique & & & & & & & & \\
    \hline

    % --- Communiquer ---
    Communiquer & Exécution des signes conventionnels & & & & & & & & \\
    \hline

    % --- Retourner en surface ---
    \multirow{5}{=}{Retourner en surface} & Maîtrise de la vitesse de remontée & & & & & & & & \\ \cline{2-10}
    & Tenue d'un palier & & & & & & & & \\ \cline{2-10}
    & Tour d'horizon & & & & & & & & \\ \cline{2-10}
    & Gonflage du gilet en surface & & & & & & & & \\ \cline{2-10}
    & Remontée en expiration contrôlée & & & & & & & & \\
    \hline

    % --- Évoluer en sécurité ---
    \multirow{2}{=}{Évoluer en sécurité} & Application des procédures mises en œuvre par le GP & & & & & & & & \\ \cline{2-10}
    & Intervention en relais & & & & & & & & \\
    \hline

  \end{longtable}
\end{center}

L'évaluation des compétences se fait en \textbf{contrôle continu}. Le niveau est validé
lorsque toutes les compétences sont acquises. Les connaissances théoriques sont évaluées à
l'oral, lors des mises en situations pratiques. Il n'y a pas d'examen écrit. L'accent est mis
sur la prévention.

\vspace{0.5cm}
\textbf{Validé le :} \makebox[6cm]{\hrulefill} \hfill \\
\\
\textbf{Nom et signature du moniteur :} \makebox[6cm]{\hrulefill}
\vspace{0.5cm}

{\small\itshape Cette fiche n'a aucun caractère obligatoire, elle est modulable et il
appartient au moniteur de l'adapter pour sa propre utilisation.}

% Prépare la numérotation pour le prochain \input{table-d-evaluation} (S9 à S16).
\addtocounter{seanceoffset}{8} (S9 à S16).
\addtocounter{seanceoffset}{8} (S9 à S16).
\addtocounter{seanceoffset}{8} (S9 à S16).
\addtocounter{seanceoffset}{8}